%%%%%%%%%%  scan idx 26  |  folio 13  %%%%%%%%%%
alg\'ebrique clos de caract\'eristique $p>0$, et que $X$ est un sch\'ema alg\'ebrique propre et lisse de dimension $3$ disons. Voir commentaires dans Exp.\ XIV n\textsuperscript{o}~2.

\subsection*{5.}
Sous la forme pr\'esent\'ee dans le n\textsuperscript{o}~4, la formule de Riemann--Roch (1.2) a \'egalement un sens pour un morphisme propre localement d'intersection compl\`ete d'espaces analytiques complexes (ou d'espaces rigide-analytiques au sens de Tate). Mais bien s\^ur la d\'emonstration donn\'ee dans [RRR] ne s'applique \`a ce cas que lorsqu'on suppose de plus que le morphisme $f$ est \underline{projectif}, et comme dans le cas alg\'ebrique, il semble qu'il faille une id\'ee essentiellement nouvelle pour traiter le cas g\'en\'eral. La formule que nous envisageons ici est une formule \`a valeurs dans un «~anneau de Chow analytique~» $\Gr^\bullet(X)$ ou plut\^ot $\Gr^\bullet(X)_{\mathbb Q}$, construit suivant le principe expliqu\'e dans 4.2 (la difficult\'e qui y est rencontr\'ee s'\'evanouissant d'ailleurs du fait qu'on est en caract\'eristique nulle, comme on le signale dans Exp.\ XIV n\textsuperscript{o}~1).

Lorsque $X$ et $Y$ sont non singuliers, il est possible de donner une autre version de la formule de Riemann--Roch en proc\'edant comme suit. Soit d'abord $X$ un espace analytique complexe quelconque, $X^{\mathrm{top}}$ le m\^eme espace muni du faisceau des fonctions complexes continues. On a donc un morphisme canonique d'espaces annel\'es
\[
X^{\mathrm{top}}\longrightarrow X,
\]
qui d\'efinit, gr\^ace au caract\`ere contravariant \'evident du foncteur $K^\bullet(X)$ en l'espace annel\'e variable $X$, un homomorphisme d'anneaux
\begin{equation}\tag{5.1}
K^\bullet(X)\longrightarrow K^\bullet(X^{\mathrm{top}}).
\end{equation}
Supposons pour simplifier $X^{\mathrm{top}}$ compact~; alors, comme dans le cas (d\'ej\`a signal\'e dans 3.1) d'un sch\'ema admettant un module inversible ample, on voit (Exp.\ II) que tout complexe parfait sur $X^{\mathrm{top}}$ est globalement isomorphe \`a un complexe born\'e localement libre en chaque degr\'e, i.e. \`a un complexe de fibr\'es vectoriels complexes sur $X^{\mathrm{top}}$, et que par suite $K^\bullet(X^{\mathrm{top}})$ est isomorphe canoniquement au groupe $K(X^{\mathrm{top}})$ bien connu des topologues [1], d\'efini en termes de fibr\'es vectoriels complexes sur l'espace compact $X^{\mathrm{top}}$. On peut donc interpr\'eter l'homomorphisme (5.1) comme \'etant un homomorphisme du $K^\bullet(X)$ d\'efini en termes de la structure analytique, dans l'anneau bien connu $K(X^{\mathrm{top}})$ des topologues. Si maintenant $f:X\longrightarrow Y$ est un morphisme d'espaces analytiques non singuliers compacts, on sait lui associer gr\^ace \`a Atiyah--Hirzebruch 
